\section{Introduction}
\label{sec:intro}

E-matching is a fundamental operation of an e-graph, used in Satisfiability
Modulo Theories (SMT) solvers for quantifier instantiation and in equality
saturation engines for applying rewrite rules. Traditional e-matching algorithms
are top-down, starting with the outermost function and recursively matching
sub-patterns.

However, recent work has proposed a new technique called relational
e-matching~\cite{relationalematching}, where database query execution techniques
are used to match syntactic patterns in an e-graph. Relational e-matching has
been  in state-of-the-art equality saturation engines such as
egglog~\cite{egglog}, but until now it has not been implemented in an SMT solver.
This is because existing relational e-matching implementations cannot be used
out of the box in an SMT solver: they must support \emph{backtracking}.

In this work, we present our initial implementation of a backtrackable,
relational e-matching algorithm in Sundance~\cite{sundance}, an SMT solver currently in
development. We describe maintaining backtrackable indices using a persistent 
data-structure and implement an e-matching algorithm based on a
left-deep binary join.

We evaluate our implementation on a set of synthetic benchmarks designed to test
e-matching and show a $25 \times$ median speedup over a single top-down
e-matching algorithm. However, the story is less rosy on real workloads. On a
set of benchmarks from the program verifier Verus~\cite{verus}, the relational
e-matching algorithm provides a median slowdown of $1.77 \times$. 

There is promise in the relational approach, but it requires more work to close
the gap on real workloads. We conclude with a discussion of future work
including better backtrackable index maps, a better understanding of real
benchmarks, multi-query optimization, and heuristics inside the query
planner. We believe that relational e-matching can be a practical component of
an SMT solver.
